\section{O teste de estresse final: quanta confiança os efeitos merecem}
\label{sec:inferenciafinal}

A Seção~\ref{sec:robustez} encerrou com um aviso: todo o efeito de
gestão medido dentro de hospital repousa em cinco conversores, e
margens de erro calculadas com tão poucas unidades tratadas tendem a
ser otimistas. Esta seção fecha a pesquisa quantificando exatamente
isso, com três instrumentos usados na literatura para situações de
poucos tratados, e resolve a pendência técnica da UTI.

\textbf{Conversores de mentira.} O teste mais intuitivo é o de
permutação: sorteamos 5 hospitais falsos entre os 309 que nunca
mudaram de gestão, fingimos que eles converteram nos mesmos anos dos
conversores reais e medimos o efeito que o modelo atribuiria a essa
conversão inventada. Repetindo o sorteio 1.999 vezes, obtemos a régua
do acaso. O resultado: a queda de custo real atribuída à OSS é comum
entre os falsos conversores (49\% deles produzem efeito tão grande
quanto o real), enquanto \textbf{a queda de tempo de permanência de
10,2\% é maior que a de 85\% dos placebos} (p igual a 0,153). Um
segundo método, o bootstrap selvagem por cluster, sabidamente
conservador nesta situação, aponta na mesma direção. A tradução: o
efeito no tempo de permanência é consistente em direção e tamanho, mas
\textbf{cinco hospitais não bastam para sustentar a afirmação com o
rigor estatístico convencional}; ele deve ser tratado como evidência
sugestiva forte, não como fato demonstrado.

\textbf{Um estimador moderno para conversões em anos diferentes.} Como
as conversões ocorreram em 2019, 2023 e 2025, a comparação clássica
pode misturar hospitais já convertidos no grupo de controle. O
estimador de Callaway e Sant'Anna evita isso comparando cada coorte de
conversão apenas com os hospitais que nunca mudaram. A
Figura~\ref{fig:cstwfe} mostra que, no ano da conversão, os dois
métodos coincidem em todos os indicadores (na mortalidade, 1,298
contra 1,301 pontos percentuais de queda); nas janelas intermediárias
o método clássico exagera um pouco o efeito do tempo de permanência
(0,25 contra 0,15 em log pontos) e as janelas longas se aproximam, com
intervalos sobrepostos (na mortalidade, 2,18 pontos negativos no
método clássico contra 2,38 no alternativo; no tempo de permanência,
0,28 contra 0,30 em log pontos). Nenhuma conclusão muda; a magnitude
intermediária do tempo
de permanência era o único ponto inflado.

\begin{figure}[H]
\centering
\subfig{fig_est_07_cs_twfe_mort.png}{Mortalidade geral, em pontos
percentuais}
\subfig{fig_est_07_cs_twfe_tmp.png}{Tempo de permanência, em log pontos}
\caption{Estudo de eventos estimado pelo método clássico (TWFE, com
intervalo de confiança) e por Callaway e Sant'Anna, que compara cada
coorte de conversão só com hospitais nunca convertidos.}
\label{fig:cstwfe}
\end{figure}

\comoler{os círculos escuros repetem o estudo de eventos da
Seção~\ref{sec:efeitogestao}; os quadrados verdes são o estimador
alternativo. Quando os dois símbolos quase se tocam, como no ano zero
e na janela longa, o resultado não depende do método; a distância nos
anos intermediários do tempo de permanência mostra o único ponto em
que o método clássico exagerava.}

\textbf{O gêmeo sintético.} Para Sorocaba e Pérola Byington, os dois
conversores com janela posterior utilizável, construímos um
\emph{controle sintético}: uma combinação ponderada de hospitais nunca
convertidos que reproduz a trajetória do hospital antes da mudança de
gestão. Se a conversão teve efeito, o hospital real deve descolar do
seu gêmeo depois dela. A Figura~\ref{fig:sintetico} traz o resultado
para a mortalidade, e ele é a ducha de água fria mais importante do
documento: \textbf{em Sorocaba o hospital real fica 1,09 ponto
percentual acima do gêmeo} após a conversão (o gêmeo, que acompanhava
a subida até 2018, também cai depois, confirmando que parte da queda
era retorno natural ao patamar); \textbf{no Pérola Byington o hospital
fica 0,95 ponto abaixo do gêmeo}, na direção esperada, mas dentro do
que sorteios de placebo produzem (p igual a 0,33). No custo real, a
lacuna do Pérola é de R\$ 354 a menos por saída (p igual a 0,28). Em
nenhum caso o descolamento é raro o bastante para valer como prova.

\begin{figure}[H]
\centering
\subfig{fig_est_06_sintetico_2081695_mort_pp.png}{Conjunto Hospitalar
de Sorocaba (CNES 2081695), conversão em 2019}
\subfig{fig_est_06_sintetico_2078287_mort_pp.png}{Pérola Byington
(CNES 2078287), conversão em 2023}
\caption{Mortalidade geral: hospital real contra o seu controle
sintético, construído só com hospitais nunca convertidos.}
\label{fig:sintetico}
\end{figure}

\comoler{a linha preta é o hospital; a azul tracejada é o gêmeo
sintético, montado para imitar o hospital antes da linha vertical. O
efeito da conversão é a distância entre as duas linhas depois da linha
vertical. Em Sorocaba a linha preta termina acima da azul; no Pérola,
abaixo, mas com o gêmeo também caindo.}

\textbf{A pendência da UTI, resolvida.} A intensidade de uso de UTI
foi reestimada com a família Gama, adequada à cauda de UTIs de uso
quase nulo que o modelo em logaritmo não acomodava. O diagnóstico de
resíduos salta de 0,814 para \textbf{0,944}
(Figura~\ref{fig:qqgama}, no Apêndice~\ref{ap:diagnosticos}),
e o resultado substantivo permanece: uso de UTI \textbf{18,2\% maior
nas OSS} (p igual a 0,005), com a variante de corte mínimo de
atividade confirmando (22,3\%). A conclusão da
Seção~\ref{sec:resultados} fica, portanto, validada também neste
indicador.

\textbf{Balanço final da fase de estimação.} O quadro que resiste a
toda a bateria de testes é: as OSS operam um regime assistencial
claramente mais intenso (mais produção, mais ocupação, mais uso de
UTI), com diferenças firmes \textbf{enquanto retrato transversal das
redes}; a única evidência com pretensão causal que resiste é a do
tempo de permanência, presente em todos os métodos, e mesmo ela
permanece sugestiva com cinco conversores, não definitiva; o custo
real por saída não mostra diferença demonstrável em nenhum método; e
a queda de mortalidade após as conversões, embora visível, não
sobrevive ao confronto com o gêmeo sintético em Sorocaba e fica sem
significância no Pérola Byington. A segunda fase da estimação
(Seção~\ref{sec:fase2}) quantificou esse último ponto da forma mais
direta possível: retirando os cinco conversores da amostra, a
diferença de mortalidade entre as redes OSS e Direta simplesmente
desaparece (o coeficiente vira 0,09 positivo, estatisticamente nulo),
ou seja, \textbf{a mortalidade menor é um fenômeno da trajetória dos
hospitais que mudaram de gestão, sem contraparte na comparação entre
as redes}. Novos anos de dados, especialmente
para os três hospitais convertidos em 2025, são o caminho natural para
transformar as sugestões em conclusões.
